% !PRINTED 513
\weberpage{513}

Wir wollen jetzt unter $x_0, y_0, \dots$ irgend welche ganze rationale Zahlen verstehen. Dann ist auch
\[ \delta_0 = \alpha x_0 + \beta y_0 + \cdots \]
ein ganzes Functional, und wir wollen beweisen, \gesperrt{dass $\delta_0$ durch $\delta$ theilbar ist.}

Wir brauchen zu diesem Zwecke nur das ganze Functional
\[ \delta t - \delta_0 = \alpha (x t - x_0) + \beta (y t - y_0) \dots \]
zu betrachten, worin $t$ eine neue Variable bedeutet. Dann ist [nach (2)]:

\weq{3}{N (\delta t - \delta_0) = D E (x t - x_0, y t - y_0, \dots),}

oder, indem wir nach absteigenden Potenzen von $t$ ordnen:

\weq{4}{N (\delta t - \delta_0) = D (t^n E + t^{n-1} E_1 + t^{n-2} E_2 + \cdots),}

worin $E_1, E_2, \dots$ nach \S.\,137, 7., 8. ganze rationale Functionale sind. Setzen wir nun
\[ C_1 = \frac{E_1}{E}, \quad C_2 = \frac{E_2}{E}, \cdots, \]
so sind, da $E$ eine Einheit ist, auch $C_1, C_2, \dots$ ganze rationale Functionale, und es folgt, wenn wir noch
\[ \frac{\delta_0}{\delta} = \eta \]
setzen, aus (4)
\[ N (\delta) N (t - \eta) = D E (t^n + C_1 t^{n-1} + C_2 t^{n-2} + \cdots), \]
oder wegen (2)

\weq{5}{N (t - \eta) = t^n + C_1 t^{n-1} + C_2 t^{n-2} + \cdots}

Da diese Function nun verschwindet, wenn $t = \eta$ gesetzt wird, so folgt, dass $\eta$ ein ganzes Functional ist (\S.\,137, 2.), und damit ist bewiesen, dass $\delta_0$ durch $\delta$ theilbar ist.

Da $x_0, y_0, \dots$ beliebige ganze rationale Zahlen bedeuten können, so schliessen wir daraus, dass die Functionale $\alpha, \beta, \dots$ selbst durch $\delta$ theilbar sind, dass also $\delta$ ein gemeinsamer Theiler der Functionale $\alpha, \beta, \dots$ ist.

Andererseits ist aber auch (nach \S.\,138, 4.) $\delta$ durch jeden gemeinsamen Theiler von $\alpha, \beta, \dots$ theilbar, und $\delta$ hat also die \gesperrt{charakteristischen Eigenschaften des grössten gemeinschaftlichen Theilers von $\alpha, \beta, \dots$} Dieselben Eigenschaften kommen aber nach \S.\,138, 7. jedem mit $\delta$ associirten Functional zu, und ebenso sind auch zwei Functionale $\delta, \delta'$, die diese doppelte Eigen-
%CONT>>
